% Ad Familiares 10.6 --- headnote
% ad-familiares-10-06, 20 March 43 BC (evening), Rome

\textit{Cicero to L. Munatius Plancus, written from Rome
on the evening of 20 March 43 BC --- Perseus dateline
\textit{Scr. Romae xiii K. Apr. vesperia. 711 (43)}, and
the closing line \textit{D. xiii K. Apr.}\ in the Latin
confirms the date of dispatch. Three months after Fam.
10.5, the political weather has changed sharply: Antony
has marched on Mutina and laid siege to Decimus Brutus;
the Senate has declared a \textit{tumultus} and sent the
consuls Hirtius and Pansa north; the Mutina campaign
is opening. Plancus has sent in two pieces of news from
Transalpine Gaul: an oral report by his legate
G. Furnius declaring loyalty to the state, and a written
letter, read out in the Senate, that talks of peace ---
which, in Cicero's reading, has been drafted in
coordination with Lepidus and is in effect a brief for
negotiation with Antony.}

\textit{The letter is one of the most concentrated pieces
of political pressure in the Plancus correspondence. The
opening section catches Plancus in the contradiction
between Furnius's voice and his own letter; the second
applies the leverage of personal friendship; the third
delivers what is in effect a public lecture on the
distinction between holding the title \textit{consul} and
deserving the name \textit{consularis}, framed around the
sharp antithesis that closes the letter --- ``not merely
no dignity but the utmost deformity.'' ``Detach yourself
at long last from those with whom no judgement of your
own but the bonds of the times have linked you'' is aimed
squarely at Lepidus. The reference to Plancus's
\textit{collega} besieged at Mutina is to Decimus Brutus,
Plancus's colleague as consul-designate for 42, and the
``foulest band of brigands'' is again the standing
\textit{Philippics} epithet for Antony's circle.}
